\subsection{残る三型と再開する具体的な入口}
\label{ns24:support-agenda}\label{ns23:support-agenda}\label{ns22:support-agenda}\label{ns20:support-agenda}

再開時は最初に \nolinkurl{00_CONTINUATION_POLICY_ja.txt} を読む．
現在の高分母支持の必要候補は\emph{三型・二十六標識}である．
全候補は五頂点を使い，三点連結，最小次数三以上，八高辺以上，$b_1\ge4$ となる．
単位枝は高々二辺の matching であり，同じ core に二本が接すれば全分母一となる．
補 graph の孤立頂点を省略すると，候補は次の通りである．
\begin{center}\small
\begin{tabular}{c|c|c|c|r}
型&高辺数&次数列&低辺補 graph&標識数\\\hline
$G_{5,14}$&8&$(43333)$&$2K_2$&15\\
$G_{5,16}$&9&$(44433)$&$K_2$&10\\
$G_{5,17}$&10&$(44444)$&$\varnothing$&1
\end{tabular}
\end{center}
識別子は詰め直さない．検算は \nolinkurl{verify_support_agenda_v23.py} である．
候補数は整数接着，atlas，幾何的非同型類の存在数ではない．

\paragraph{高分母全体の原始性．}
非空 $H$ には八辺以上ある．全高分母を割る素数があるとすると，素数支持が $H$ と一致するため，
密な素数支持の定理はその素数を五，支持を $K_{2,2,1}$ に強制する．
補 graph の二辺は実際には単位分母であるが，合同分類はその和を零 modulo five にするため矛盾する．
従って高分母全体の最大公約数は一である．

二本の単位枝がある場合には任意七本の分母の最大公約数も一となり，
$G_{5,14}$ で五の倍数の分母は高々六本である．
密な $G_5=K_{2,2,1}$ と一本の実単位枝が共存する場合は $H=G_{5,16}$ に限られ，
五進支持の外のもう一本は単位枝と頂点を共有せず，分母 $q\equiv9\pmod{10}$，$q\ge9$ となる．
合同条件だけで整数 lifting を保証してはいない．

\paragraph{優先候補：単位枝の node と小分母の位置．}
Book と $G_{5,15}$ は全排除済みであり，その専用探索は再開しない．
$G_{5,14}$ の二本の単位枝には，正の三高 core 同士が接する．
各枝の node 配分は $\{5,0,0\}$ または $\{4,1,0\}$ である．
共通負位置の分母対は，五nodeなら $(3,7),(5,5),(7,3)$，
四nodeなら負 pair の $(3,3)$ または正 pair の $(3,2R+1),(2R+1,3)$ となる．
後者の $R\ge1$ は自由であり，表の有限性を全分母の有限性と解釈しない．

対向低 $D$ を相殺すると，単位枝の三つの外側 pair の少なくとも一つが正になる．
五nodeの場合は他の二つがともに正で，対応する高辺の node は零になる．
$G_{5,14}$ の五node位置はもう一方の単位枝の端点に限られ，四高 core は除かれる．
これらと三高の matching 性を同時に課し，
両単位枝の共通負位置・node 四／五・小分母三／五／七の位置を先に有限分類する．
全行和・pair products・対向高 $D$・各高辺の node 総数を復元時に戻す．
その後に第三 core の単一 pair による線形復元，非共鳴の分母比較，
共鳴族と小分母例外を組み合わせて閉路条件を解く．
固定中心四点への還元は $G_{5,15}$ 固有の構造なので流用しない．

\paragraph{並行候補：三高同時非負性の残る格子領域．}
固定 $(s,t)$ ごとの全係数復元を基礎として，$|\delta|\ge2$ では
尺度の正値性が許す $x,n$ の幅を元の pair 残余と組み合わせる．
$\delta=\pm1$ は本版までに全排除された．
比例領域 $\delta=0$ は $x=y=0,z\ge1,s<t,0\le n\le z$ の零軸だけが残る．
canonical な向きを保ち，標識交換で $s<t$ を消してはならない．
固定 $(s,z)$ で各 $n=0,\ldots,z$ の真の四次式を全整数根まで解けるうえ，
\[
 b+1\mid C_n,
 \qquad C_n=n\bigl((sz)^2-3sz+3\bigr)+(s^2-1)z-3s+1>0
\]
も必要である．$b+1\le C_n\le C_z$ を使い，
$n=0,n=z,0<n<z$ の整数根・合同障害を分ける．
零余りの二共鳴は排除済みだが，$s,z$ 自体には一様上界を得ていない．

退化面 $a=0$ では $x=0,R_0=1$ を全排除したため，固定 $s$ の完全有限復元がある．
$s=1,2$ は解がなく，$s\ge3$ では
\[
 0\le x\le s-1,\qquad 1\le R_0\le(s-x)(s-1)+1
\]
を取り，$x=0,R_0=1$ を除く．
$x\ge1$ なら $b\le\lfloor(s-1)/x\rfloor$，$x=0$ なら
$b\le1+s^2(s-1)^2$ である．二次式で $F$ を，元残余の一次式で $y$ を復元する．
退化面全体の不存在や $s$ の一様上界は未証明である．
正側軸 $\delta=1,x=0$ と退化例外 $x=0,R_0=1$ は未解決リストから外す．
反例が現れれば，全原方程式・奇数性・六符号・正 weights を確認して構成へ転換する．

\paragraph{四高の二符号型と幾何への転換．}
$G_{5,16}$ の唯一の単位枝にも上の三角形表を適用する．
$G_{5,16},G_{5,17}$ の四高 core では正符号 graph は $2K_2$ または $P_4$ に限られる．
どちらも正 weights を持つため，次は元の整数 pair・行和条件を扱う．
$P_4$ の一般排除または具体的構成を独立課題とし，三高 matching 定理を流用しない．
$K_4$ を部分 graph として含むことだけでも支持を排除できない．
整数候補が残れば全 core の weights，basin areas，integer framings，作用切片へ進む．
五次・六次 atlas の同値，不変量，別の実際の退化，open-corridor surgery への転換も，
研究全体の面白さとインパクトに応じて判断する．
